\section{A bounded backward intertwiner}\label{sec:intertwiner}
We seek an operator which sends $e_j$ in the fiber at $x$ to a positive multiple of $e_{j-1}$ in the fiber at $\theta x$. Commutation with the forward shift determines a recurrence for those multiples. The opposite logarithmic drifts at the two ends of the fiber make its solutions bounded after normalization.

\subsection{Two-ended drift}
Write $\Delta a_j(x)=a_j(\theta x)-a_j(x)$. A block $I_m$ has an exceptional residue at $x$ when $T_m(x)=N_m-1$. Since
\[
 \sum_{m\ge1}\mu\{T_m=N_m-1\}=\sum_{m\ge1}2^{-m}<\infty,
\]
the Borel--Cantelli Lemma shows that almost every $x$ has only finitely many such blocks. Intersect this conull set over all iterates of $\theta$ to obtain a measurable $\theta$-invariant conull set $X_*$. No independence of the exceptional events is needed.

\begin{lemma}\label{lem:drift}
For every $x\in X_*$ there is a finite $C(x)$ such that, for all $n,q\ge1$,
\begin{align}
 \sum_{k=0}^{n-1}\Delta a_k(x)&\ge Dn-C(x),\label{eq:positive-drift}\\
 \sum_{k=-q}^{-1}\Delta a_k(x)&\le-Dq+C(x).\label{eq:negative-drift}
\end{align}
\end{lemma}
\begin{proof}
For $k\in I_m$, Equation~\eqref{eq:residue-step} gives
\[
 \Delta a_k^+(x)=D-DN_m\mathbf1_{\{T_m=N_m-1\}}(x),
 \qquad \Delta a_k^-(x)=-\Delta a_k^+(x).
\]
Both central increments $\Delta a_0^\pm$ vanish. Thus the first sum is
\[
 D(n-1)-D\sum_{\substack{m\ge1\\T_m(x)=N_m-1}}
       N_m\bigl|I_m\cap[1,n)\bigr|.
\]
Under $k\mapsto-k-1$, the negative-index sum is the negative of the same expression with $n=q$. Both claims follow on taking
\[
 C(x)=D+D\sum_{\substack{m\ge1\\T_m(x)=N_m-1}}N_m|I_m|<\infty.
\]
\end{proof}

The constant $C(x)$ may be unbounded and nonintegrable as a function of $x$. Only its pointwise finiteness is used below.

\subsection{Solving and normalizing the recurrence}
For positive coefficients $b_j(x)$, sending $e_j$ at $x$ to $b_j(x)e_{j-1}$ at $\theta x$ would commute with $S$ precisely when
\[
 b_j(x)\beta_{j-1}(\theta x)=\beta_j(x)b_{j+1}(x).
\]
Figure~\ref{fig:square} shows the two paths in this coefficient identity. The left side is obtained by moving backward and then forward, and the right side in the opposite order.

\begin{figure}[ht]
\centering
\begin{tikzpicture}[>=Stealth, every node/.style={font=\small}]
\node (a) at (0,2) {$e_j\text{ at }x$};
\node (b) at (6,2) {$e_{j+1}\text{ at }x$};
\node (c) at (0,0) {$e_{j-1}\text{ at }\theta x$};
\node (d) at (6,0) {$e_j\text{ at }\theta x$};
\draw[->] (a)--node[above] {$S_x:\ \beta_j(x)$} (b);
\draw[->] (c)--node[below] {$S_{\theta x}:\ \beta_{j-1}(\theta x)$} (d);
\draw[->] (a)--node[left] {$V_x:\ b_j(x)$} (c);
\draw[->] (b)--node[right] {$V_x:\ b_{j+1}(x)$} (d);
\end{tikzpicture}

\caption{The coefficient identity at a fixed input fiber. Both paths take $e_j$ at $x$ to a multiple of $e_j$ at $\theta x$. The horizontal maps move forward in the coordinate, while the vertical maps move backward and translate the base.}\label{fig:square}
\end{figure}
 We first solve this equation with initial value one, then divide all coefficients in each fiber by a common bound.
For every $x\in X$ set $\widetilde b_0(x)=1$ and define positive coefficients at all integer indices by
\begin{equation}\label{eq:b-recursion}
 \widetilde b_{j+1}(x)
 =\widetilde b_j(x)\frac{\beta_{j-1}(\theta x)}{\beta_j(x)}
 \qquad(j\in\Z).
\end{equation}
Every coefficient is a finite product of positive measurable functions and their reciprocals, so it is finite and measurable.

\begin{lemma}\label{lem:normalization}
For $x\in X_*$, the coefficients $\widetilde b_j(x)$ tend to zero as $j\to+\infty$ and as $j\to-\infty$. Consequently there are measurable functions $b_j:X\to[0,1]$ which are positive on $X_*$ and satisfy
\begin{equation}\label{eq:commuting-coefficients}
 b_{j+1}(x)\beta_j(x)=b_j(x)\beta_{j-1}(\theta x)
 \qquad(x\in X,\ j\in\Z).
\end{equation}
\end{lemma}
\begin{proof}
Taking logarithms in \eqref{eq:b-recursion} and summing forward gives
\begin{align}
 \log\widetilde b_n(x)
 &=\sum_{k=0}^{n-1}\bigl(a_k(x)-a_{k-1}(\theta x)\bigr)\notag\\
 &=-\sum_{k=0}^{n-1}\Delta a_k(x)
       +a_{n-1}(\theta x)-a_{-1}(\theta x),\qquad n\ge1.\label{eq:b-positive}
\end{align}
Summing backward from index $0$ gives the different sign
\begin{align}
 \log\widetilde b_{-q}(x)
 &=-\sum_{k=-q}^{-1}\bigl(a_k(x)-a_{k-1}(\theta x)\bigr)\notag\\
 &=\sum_{k=-q}^{-1}\Delta a_k(x)
       +a_{-q-1}(\theta x)-a_{-1}(\theta x),\qquad q\ge1.\label{eq:b-negative}
\end{align}
Here $a_{-1}=c$. The boundary indices have different forms:
$a_{n-1}=a_{n-1}^+$, whereas $a_{-q-1}=a_q^-$. Using $a_0^\pm=c$ and Equation~\eqref{eq:a-error}, we obtain
\[
 a_{n-1}(\theta x)-c\le c(n-1)+4D(n-1)^{1/3},\qquad
 a_{-q-1}(\theta x)-c\le cq+4Dq^{1/3}.
\]
Combining these bounds with Lemma~\ref{lem:drift}, we obtain the explicit estimates
\begin{align*}
 \log\widetilde b_n(x)&\le -(D-c)n+C(x)-c+4D(n-1)^{1/3},\\
 \log\widetilde b_{-q}(x)&\le -(D-c)q+C(x)+4Dq^{1/3}.
\end{align*}
Since $D-c=90>0$, both right sides tend to $-\infty$. These estimates are pointwise in $x$; they do not require an integrable bound on $C(x)$.

The countable supremum
\[
 L(x)=\sup_{j\in\Z}\widetilde b_j(x)
\]
is measurable, and $1\le L(x)<\infty$ on $X_*$. Define
\[
 b_j(x)=\begin{cases}\widetilde b_j(x)/L(x),&x\in X_*,\\0,&x\notin X_*.
 \end{cases}
\]
All indices at a fixed $x$ are divided by the same number, so \eqref{eq:b-recursion} yields \eqref{eq:commuting-coefficients}. Off $X_*$ both sides of that identity vanish.
\end{proof}

Define a backward shift $V_x:K\to K$ by $V_xe_j=b_j(x)e_{j-1}$. Every $V_x$ has norm at most one, and $x\mapsto V_xv$ is measurable for every $v\in K$ by finite-coordinate approximation. It defines $V\in\B(H_0)$ by
\begin{equation}\label{eq:V-definition}
 (Vf)(\theta x)=V_xf(x),
 \qquad\text{equivalently }(Vf)(\theta x,j-1)=b_j(x)f(x,j).
\end{equation}
Invariance of $\mu$ gives $\norm{Vf}^2=\int_X\norm{V_xf(x)}^2\,d\mu(x)\le\norm{f}^2$.

\begin{proposition}\label{prop:V}
The contraction $V$ commutes with $S$. For every $x\in X_*$ and $h\in\Z$,
\[
 \overline{V_x\Ktail{h}}=\Ktail{h-1},\qquad
 \overline{V_xK}=K.
\]
\end{proposition}
\begin{proof}
At the same output base point $\theta x$, the two compositions have coefficients
\[
 S_{\theta x}V_xe_j=b_j(x)\beta_{j-1}(\theta x)e_j,
 \qquad
 V_xS_xe_j=\beta_j(x)b_{j+1}(x)e_j.
\]
They agree by \eqref{eq:commuting-coefficients}, first on the basis and then on $K$. Equation~\eqref{eq:V-definition} now gives $SV=VS$. On $X_*$ every $b_j(x)>0$. The image of a tail therefore contains every basis vector in the preceding tail, up to a nonzero scalar, and is contained in that preceding tail. This proves the density assertions.
\end{proof}


The density conclusion, rather than surjectivity, is exactly what is needed for closed invariant subspaces. No bounded inverse of $V_x$ or $V$ is asserted.
